\chapter{Recurso hídrico, estiaje y energía}
\section{Datos públicos recuperados y su naturaleza}
La estación hidrológica Limbani del mapa público del SENAMHI tiene código 260501, categoría HLG y condición de transmisión diferida. El visor permite identificar la estación, pero la descarga de datos está protegida por una verificación de usuario. No se obtuvo una serie de caudales ni su curva nivel--caudal. El boletín de diciembre de 2025 consultado informa \textbf{nivel en metros}; convertirlo en caudal sin una curva de gasto sería incorrecto. Tampoco se debe confundir esta estación con la meteorológica CO 114049. Sus coordenadas publicadas y los inventarios históricos no coinciden exactamente y se registran por separado (H03 y H17).

Se recuperaron íntegramente los productos públicos PISCO de SENAMHI, publicados en HydroShare: ARNOVIC v1.1 diario y GR2M v2.0 mensual (H01--H02). Ambos abarcan 1981--2025. Son \textbf{caudales simulados de modelos hidrológicos}, no aforos en Limbani. La extracción del COMID 9080887 contiene 16\,436 días consecutivos y 540 meses, en \(\mathrm{m^3/s}\), sin valores faltantes ni negativos. Los archivos NetCDF, sus vectores y sus huellas SHA-256 se conservan para reproducir el cálculo. La licencia publicada es CC BY 4.0; los autores y la procedencia se mantienen en el registro de fuentes.

El tramo cartográfico pasa a unos 25 m de la referencia CM-B, pero su caudal representa la \textbf{salida de una cuenca aguas abajo}. La casa de máquinas está aproximadamente al 40\% del recorrido del tramo de 32.6 km. Aplicar directamente ese caudal a la captación sobreestimaría el recurso. Por ello se estudia una transposición espacial explícita, con incertidumbre.

\section{Nombre de los ríos y antecedentes próximos}
La denominación río Limbani está respaldada por SENAMHI y los antecedentes locales. No se ha verificado que ``río Para'' sea el nombre oficial de un río aprovechable en el emplazamiento. Se investigó Phara, distrito vecino, sin declarar que Para y Phara sean equivalentes. La resolución ANA 370-2017 identifica una captación en el río \textbf{Cullucachi}, y otro antecedente ANA menciona el río Comunidad en Phara (H07 y H18). La licencia de la minicentral de Phara concede 90 l/s y 2\,838\,240 \(\mathrm{m^3/a}\) para 100 kW; se trata de un derecho de uso de \emph{otra} central, no de un registro natural de caudales ni de una autorización para Limbani.

La tesis de Goyzueta Arce (2017), de la Universidad Nacional del Altiplano, aporta un diseño de Phara con turbina Pelton, \(Q=0.24\ \mathrm{m^3/s}\) y salto neto de 134.25 m. Su cuadro de cinco años no identifica los años calendario. No se adopta como serie temporal oficial ni se transfiere mediante regresión: faltan aforos concurrentes, áreas verificadas, precipitación comparable y control de usos del agua. Su mayor salto tampoco puede trasladarse a Limbani.

\section{Crecidas y seguridad: análisis independiente}
El máximo diario simulado, \dato{Qhigh} \(\mathrm{m^3/s}\), no es pico instantáneo ni avenida de período de retorno conocido. Para bocatoma y casa se requieren extremos, secciones, niveles, socavación y sedimentos, con períodos de retorno justificados por riesgo. Este estudio de seguridad se desarrolla separado de generación; el máximo simulado no certifica las obras.

\section{Transposición por área y comprobaciones cartográficas}
Se delimitó una cuenca académica para la captación propuesta mediante un modelo de elevación público Mapzen, con relleno de depresiones y dirección de drenaje D8. El algoritmo conserva la acumulación de celdas y se verificó con pruebas independientes. La posición se ajusta a la celda de mayor acumulación dentro de 250 m: esta tolerancia reconoce que la captación todavía no está levantada en campo.

El área del atributo del vector PISCO es \(\dato{Amerc}\ \mathrm{km^2}\), pero coincide con el área calculada en \textbf{Web Mercator}. Esa proyección distorsiona las superficies. Se reemplaza por el área geodésica WGS84, \(A_r=\dato{Aref}\ \mathrm{km^2}\). La cuenca estimada en captación tiene \(A_c=\dato{Acapt}\ \mathrm{km^2}\). El factor de estudio resulta:
\begin{equation}
 Q_c(t)=\alpha Q_r(t),\qquad \alpha=\frac{A_c}{A_r}=\dato{Alpha}.
\end{equation}
Esta relación supone rendimientos específicos uniformes entre una subcuenca y su cuenca de referencia. No es una regresión ajustada ni una regionalización calibrada. Cerca del \(\dato{NestedPct}\%\) de la subcuenca delimitada está dentro del polígono PISCO, y sus áreas cambian menos del 1\% al repetir la delimitación a resoluciones de 30 y 90 m. Esto verifica estabilidad geométrica local, \textbf{no exactitud del caudal}.

Existen dos limitaciones relevantes. Primero, la cota media de la subcuenca es aproximadamente \dato{ElevationLocal} m, frente a \dato{ElevationRef} m en el polígono de referencia; la precipitación, evapotranspiración y almacenamiento pueden diferir. Segundo, el ajuste del DEM en la salida PISCO captura una confluencia ambigua y no reproduce satisfactoriamente el área de referencia. Ese control \textbf{no fue aceptado}. No se normaliza el caudal por una cuenca de salida D8 arbitraria. Se mantiene el área geodésica del producto y se examinan sesgos de 0.5, 0.7, 1.0 y 1.3; no son intervalos de confianza estadísticos.
\begin{figure}[htbp]\centering\includegraphics[width=.95\linewidth]{figuras/cuenca-r03.png}\caption{Cuenca propuesta y referencia PISCO. Delimitación de gabinete; no es una cuenca aforada ni un levantamiento de campo.}\end{figure}

\section{Estacionalidad y curva de duración}
\input{generado/hidro_estadisticas}
El caudal diario medio transpuesto es \dato{Qmean} \(\mathrm{m^3/s}\), y su Q95 diario es \dato{Qlow} \(\mathrm{m^3/s}\). Q95 significa excedencia aproximada durante el 95\% de los días simulados del período; no garantiza ese caudal cada año ni autoriza una extracción. La curva se construye ordenando los valores de mayor a menor y asignando posición de Weibull \(p=m/(n+1)\). Las cuantificaciones del cuadro se interpolan sobre la distribución empírica y conservan resolución diaria.
\begin{figure}[htbp]\centering\includegraphics[width=.94\linewidth]{figuras/caudales-r03.png}\caption{Serie transpuesta y variación mensual. Las dos líneas proceden de modelos públicos, no de mediciones en la captación.}\end{figure}
\begin{figure}[htbp]\centering\includegraphics[width=.90\linewidth]{figuras/duracion-r03.png}\caption{Duración diaria del caudal natural estimado y del caudal utilizable bajo reservas de estudio.}\end{figure}
La media mensual de GR2M difiere en \dato{MonthlyBias}\% de la media diaria agregada a meses de ARNOVIC; la correlación mensual es aproximadamente 0.76. Los modelos comparten información hidrometeorológica y no constituyen validación independiente. La suavización mensual oculta los días críticos. Su potencia con excedencia mensual de 95\%, \dato{PmonthlyLow} kW en usuarios, no reemplaza al indicador diario ni demuestra continuidad del servicio.

\section{Reserva ecológica, otros usos y balance de agua}
La RJ 267-2019-ANA distingue métodos y alcances para definir el caudal ecológico según la categoría ambiental y la coordinación con la AAA. No establece un porcentaje universal del 15\%. Para este ejercicio se reserva el 15\% de la media climatológica de cada mes, \textbf{como supuesto de sensibilidad}, hasta disponer del estudio hidroecológico y de la decisión de ANA. Se deja además una reserva ilustrativa de 0.005 \(\mathrm{m^3/s}\) para otros usos; tampoco es un inventario de derechos.
\begin{align}
 Q_e(t)&=\min[Q_c(t),\,0.15\overline{Q}_{c,mes}],\\
 Q_o(t)&=\min[0.005,\,\max(0,Q_c-Q_e)],\\
 Q_a(t)&=\max(0,Q_c-Q_e-Q_o),\\
 Q_c(t)&=Q_e(t)+Q_o(t)+Q_{t,\,posible}(t)+Q_{remanente}(t).
\end{align}
El caudal posible de turbinar se limita por equipos y caudal máximo de diseño. La serie exportada describe \textbf{capacidad potencial}, no una consigna real de extracción horaria. En despacho parcial el caudal realmente turbinado sería menor, y el excedente debe quedar en el río o descargarse por dispositivos dimensionados. La reserva ecológica se mantiene en el tramo derivado, sin sustituirla por la restitución aguas abajo.

La resolución ANA de 2019 para uso poblacional en Limbani reconoce 3.7 l/s de cuatro manantiales. No corresponde restar automáticamente ese volumen del río en esta captación: deben identificarse conexiones hidráulicas, titularidad y derechos concurrentes mediante el RADA y el análisis de disponibilidad hídrica. No se acreditan disponibilidad ni licencia obtenidas.

\section{Potencia y energía: comparación de alternativas}
Se utiliza \(P=\rho gQH_n\eta_t\eta_m\eta_g\), con pérdidas de conducción, eficiencias dependientes de carga, auxiliares y pérdidas eléctricas calculadas por separado. Las curvas parciales adoptadas son hipótesis académicas; se requieren curvas del proveedor para seleccionar equipos. La disponibilidad mecánica de 95\% se aplica a la \textbf{energía esperada}, sin convertirla en una garantía de potencia firme. La demanda se mantiene con el perfil 2026 para someterla a los 45 años climáticos simulados. Cada año se normaliza a 365 días.
\input{generado/alternativas}
\input{generado/hidro_meses}
El caso de 600 kW permite aproximadamente \dato{Pusers} kW en usuarios en régimen nominal, frente a \dato{Demand} kW de demanda máxima estimada. Sin embargo, el indicador diario de excedencia de 95\% es solamente \dato{Pfirmstudy} kW y el mínimo del modelo es \dato{PMin} kW. Hay \dato{DaysZero} días equivalentes por año sin generación bajo las reservas adoptadas. \textbf{Ni 150 ni 200 kW pueden ofrecerse como potencia continua en estiaje}. En julio, la media de potencia disponible es \dato{PJulMean} kW y el indicador diario de excedencia de 95\% es \dato{PJulLow} kW.

La capacidad hidroeléctrica anual disponible es \dato{Epotential} MWh, pero la energía que coincide con las cargas es \dato{EHidro} MWh. No se vende ni valora como ingreso la energía sobrante. El déficit hidroeléctrico es \dato{Edef} MWh/año; requiere respaldo o reducción de servicio. Incrementar de 400 a 600 kW cubre mejor el pico, aunque su ganancia de energía útil es pequeña frente al costo del tercer grupo. El escenario de 600 kW es una \textbf{envolvente conservadora de cobertura}, no la demostración de un óptimo económico.
\input{generado/sesgo}
\begin{figure}[htbp]\centering\includegraphics[width=.94\linewidth]{figuras/oferta-demanda-r03.png}\caption{Demanda sintética, capacidad variable y déficit de energía. No es un registro medido de operación.}\end{figure}
